\begin{afterword}
\summaryhead

The post starts from an invented case. Sandy is obese. Her husband blames her, her doctor prescribes drugs and surgery, and her sister calls obesity a valid lifestyle. Similar disputes, the author says, surround alcoholism, depression, attention deficit disorder and other ``marginal conditions.''

Drawing on earlier LessWrong posts, he treats ``disease'' as a cluster of six features: biological cause, involuntariness, rarity, unpleasantness, discreteness and medical treatment. Cancer has all six. Obesity has some. Asking whether it is ``really'' a disease is then empty, except that the label decides two further questions: whether the person gets sympathy, and whether they may seek medical treatment.

He answers these two questions directly from a ``determinist consequentialist'' position. Everything is biological, but some conditions respond to social pressure and some do not. One should condemn when condemnation reduces a condition enough to be worth the hurt feelings, and sympathize otherwise. People should be free to use any effective treatment. He attributes opposition to medical fixes to status quo bias and to moral signaling.

\responsehead

The title promises to dissolve a question. The post does not dissolve it. It moves the question from medicine into ethics, and there it picks a side. It admits that the side it picks is ``philosophy that the majority of the human race would disavow.'' It offers the reader a place in the minority, ``We believe it's biology all the way down,'' and a vocabulary for dismissing those outside it: they show status quo bias, or they are protecting their moral standing.

The answer it reaches is not independent of the disease concept, as the post promised it would be. Its test for sympathy is whether a condition is ``completely immune'' to social influence. That is the second of its own six criteria, ``completely immune to the operations of free will,'' with one phrase exchanged for another. Its test for treatment, ``available and effective,'' is close to the sixth. The disease node has been removed, and one of the criteria that fed it is now doing its work.

The rule for condemnation is also an empirical claim, and the post brings no evidence to it. Whether shaming reduces a condition is a question of fact. The post settles it for laziness by assertion and leaves it open for obesity, the case it began with. A review published the following month found that the evidence on obesity pointed against the husband. Because the rule counts incidence across a population, it also permits condemning people who cannot change, provided others are deterred. The view of blame is Schlick's, from 1931. The post does not answer the objection Campbell raised against it in 1951. Nor does it answer the epigraph. Treating condemnation as a dose fits the worry about ``personhood'' rather than meeting it (this is an inference). The observation that the disease question often stands in for questions about blame and treatment is useful.

The details are loose throughout, in a post that offers to replace a muddled debate with clear thinking. The metals in the borrowed example are swapped. Diabetes is said to meet every criterion and fails two. The commonest form of dwarfism is treated as the tail of a curve. The evidence of social attitudes is a satirical headline. The reversal test is misapplied. The signaling pattern it borrows is about problems that ``don't threaten them personally,'' and it is applied to a problem that does.

In short: it promises to dissolve a question and instead answers it from a theory it admits most people reject; the answer is one of its own criteria renamed, and its rule for condemning people rests on a question of fact it never investigates.
\end{afterword}
